% This template has been taken and adapted from ICML 2021 submission file available under Creative Commons CC BY 4.0.

% The main.tex ICML 2021 submission file contains the following attribution:

% "This document was modified from the file originally made available by
% Pat Langley and Andrea Danyluk for ICML-2K. This version was created
% by Iain Murray in 2018, and modified by Alexandre Bouchard in
% 2019 and 2021. Previous contributors include Dan Roy, Lise Getoor and Tobias
% Scheffer, which was slightly modified from the 2010 version by
% Thorsten Joachims & Johannes Fuernkranz, slightly modified from the
% 2009 version by Kiri Wagstaff and Sam Roweis's 2008 version, which is
% slightly modified from Prasad Tadepalli's 2007 version which is a
% lightly changed version of the previous year's version by Andrew
% Moore, which was in turn edited from those of Kristian Kersting and
% Codrina Lauth. Alex Smola contributed to the algorithmic style files."

% Link to the original template: https://icml.cc/Conferences/2021/StyleAuthorInstructions

% The adaptation of the ICML 2021 submission file was done at the EPFL and Idiap Research Institute by Alina Elena Baia, Darya Baranouskaya and Olena Hrynenko (equal contribution).

\documentclass{article}

% importing the style files
\input{source/preamble}
\usepackage[accepted]{source/icml2021}
\usepackage{subcaption}
\usepackage{float}
\usepackage{multirow}


%%%%%%%%% ADD YOUR GROUP NUMBER %%%%%%%%%
\icmltitlerunning{EE-559 Deep Learning: Group Mini-Project, Group 20}

\begin{document}

\twocolumn[
%%%%%%%%% ADD TITLE OF YOUR GROUP MINI-PROJECT %%%%%%%%%
\icmltitle{Implicit Hate Speech Against Swiss Minorities}

%%%%%%%%% ADD NAMES OF THE GROUP MEMBERS %%%%%%%%%
\begin{icmlauthorlist}
\icmlauthor{Grange Louis}{ch} \quad
\icmlauthor{Oussama Yazidi}{ch} \quad
\icmlauthor{Joshua Oyewole Oyebanji}{ch} \\
\end{icmlauthorlist}

% adding affiliation information
\icmlaffiliation{ch}{}
\vskip 0.05in
\centering
%%%%%%%%% ADD YOUR GROUP NUMBER %%%%%%%%%
Group 20
\vskip 0.3in
]

% command for printing affiliation
\printAffiliationsAndNotice{}

%%%%%%%%% ADD ABSTRACT %%%%%%%%%
\begin{abstract}
Hate speech detection systems are predominantly trained on US-centric datasets, leaving regional minorities from other countries underrepresented. We extend the ToxiGen benchmark to three Swiss-specific groups (cross-border workers, Portuguese residents, asylum seekers), generate synthetic samples with \texttt{gpt-4o-mini}, and fine-tune \texttt{tomh/toxigen\_roberta} on the resulting dataset, ToxiGen+. The baseline achieves only 26\% hate-class recall on Swiss content; fine-tuning on ToxiGen+ raises macro F1 to 0.998 with near-perfect hate recall.

%%%%%%%%% ADD KEYWORDS %%%%%%%%%
% Make sure to keep keywords on a new line.
\textbf{Keywords:}

\end{abstract}
%%%%%%%%% ADD INTRODUCTION SECTION %%%%%%%%%
\section{Introduction}
\label{sec:intro}

Automated hate speech detection is a key component of content moderation, yet most systems are trained on English-language, US-centric datasets (covering groups such as Mexican, Black, or Asian Americans) that generalize poorly to other regions and communities.

Switzerland presents a compelling case: cross-border workers (\textit{frontaliers}), Portuguese residents, and asylum seekers are frequent targets of online hostility, yet no dedicated dataset exists for these groups. We address this gap by constructing ToxiGen+ \cite{hartvigsen2022toxigen}, augmenting ToxiGen with Swiss-specific synthetic samples generated by \texttt{gpt-4o-mini}, and fine-tuning \texttt{tomh/toxigen\_roberta} on the result. Our contributions are: (i) a 3,000-sample Swiss dataset balanced between hate and neutral content, (ii) a fine-tuned classifier adapted to the Swiss context, and (iii) a comparative evaluation showing the limits of existing models on regional minorities.



%%%%%%%%% ADD RELATED WORK SECTION %%%%%%%%%
\section{Related Work}
\label{sec:related_work}

Early work on hate speech detection relied on crowd-sourced Twitter annotations, limited to explicit slurs and a narrow set of US-centric identity groups. The introduction of transformer-based models and RoBERTa \cite{liu2019roberta} shifted the field towards fine-tuning approaches, substantially improving detection of nuanced hate speech. HateBERT \cite{caselli2021hatebert} further showed that domain-adaptive pre-training on abusive content benefits downstream classifiers.

ToxiGen \cite{hartvigsen2022toxigen} addressed implicit hate speech at scale by using GPT-3 to generate 274k statements targeting 13 minority groups. While a significant step forward, its coverage remains limited to groups prominent in the American context, leaving Swiss and European minorities entirely absent. Our work fills this gap by augmenting ToxiGen with Swiss-specific groups.

%%%%%%%%% ADD METHOD SECTION %%%%%%%%%
\section{Method}
\label{sec:method}

\subsection{Base model and problem framing}
We build on \texttt{tomh/toxigen\_roberta}, a publicly released RoBERTa-large checkpoint (24 layers, 1024 hidden dimensions, $\approx$355\,M parameters) fine-tuned on the original ToxiGen corpus for binary toxicity classification \cite{liu2019roberta, hartvigsen2022toxigen}. Our contribution is framed as \emph{second-stage} fine-tuning: rather than training a classifier from a generic encoder, we adapt an already toxicity-specialised model to a new domain (Swiss-context implicit hate) while preserving its coverage of the 13 identity groups already represented in ToxiGen.

Formally, the model maps an input sequence $x$ to a label $\hat{y} \in \{0, 1\}$ via
\begin{equation}
  \hat{y} = \arg\max_{c} \; W \, h_{\texttt{[CLS]}}(x) + b,
  \label{eq:argmax}
\end{equation}
where $h_{\texttt{[CLS]}}(x) \in \mathbb{R}^{1024}$ is the \texttt{[CLS]} representation produced by the RoBERTa encoder and $W \in \mathbb{R}^{2 \times 1024}$ is the classification head.

\subsection{Dataset construction}
We built two datasets for our experiments. The first, ToxiGen-16k, consists of 16,000 samples drawn uniformly at random from the original ToxiGen dataset, covering its 13 identity groups, and serves as our baseline.

The second, ToxiGen+, extends ToxiGen with three Swiss-specific minority groups: cross-border workers, Portuguese residents, and asylum seekers. For each group, hate and neutral sentences were manually written as prompts and fed to \texttt{gpt-4o-mini} to generate 1,000 samples per group (500 hate, 500 neutral). These 3,000 Swiss samples were merged with 13,000 ToxiGen samples drawn uniformly at random, yielding a balanced dataset of approximately 16,000 entries. All samples share the same schema: \texttt{text}, \texttt{labels} ($1 =$ hate, $0 =$ neutral), and \texttt{target\_group}. Table~\ref{tab:datasets} summarises the composition of both datasets.

\begin{table}[H]
\centering
\caption{Dataset composition.}
\label{tab:datasets}
\small
\begin{tabular}{lrrr}
\hline
\textbf{Dataset} & \textbf{Total} & \textbf{Hate} & \textbf{Neutral} \\
\hline
ToxiGen-16k & 16{,}000 & 8{,}000 & 8{,}000 \\
ToxiGen+ & 16{,}000 & 8{,}000 & 8{,}000 \\
\quad of which Swiss & 3{,}000 & 1{,}500 & 1{,}500 \\
\hline
\end{tabular}
\end{table}

\subsection{Fine-tuning procedure}
We fine-tune the entire $\approx$355\,M-parameter network end-to-end, attaching a two-class classification head to the \texttt{[CLS]} token and minimising the standard cross-entropy loss
\begin{equation}
  \mathcal{L} = -\frac{1}{N}\sum_{i=1}^{N} y_i \log p_i + (1-y_i)\log(1-p_i),
  \label{eq:ce}
\end{equation}
where $p_i$ is the predicted probability of the hate class for example $i$. Training is handled via the HuggingFace \texttt{Trainer} API. Each corpus is partitioned 90/10 into train and validation splits by a label-stratified procedure with fixed seed 1, and no class weighting or resampling is applied. Optimisation uses AdamW with learning rate $2\times10^{-5}$, batch size 16, weight decay 0.01, warmup ratio 0.06, fp16 mixed precision, and a maximum sequence length of 256 tokens. We train for up to five epochs with early stopping (patience 2) on validation macro F1, retaining the best checkpoint. The same recipe is applied to both corpora, so the resulting models (ToxiGen and ToxiGen+) differ only in their training data.


%%%%%%%%% ADD VALIDATION SECTION %%%%%%%%%
\section{Validation}
\label{sec:validation}

\subsection{Evaluation protocol}
Predictions are obtained by applying Eq.~\ref{eq:argmax} to the two-class logits, which corresponds to an implicit decision threshold of 0.5. We report accuracy, macro F1 (the unweighted mean of the per-class F1 scores), and hate-class recall. Macro F1 is adopted as the primary metric because, by weighting both classes equally, it penalises a model that maximises accuracy by collapsing onto the majority class, which is precisely the failure mode that matters for a hate-speech detector. Per-target-group macro F1 additionally serves as a fairness lens across identity groups.

\subsection{Leakage audit and paired benchmarks}
Because the two models are trained on different corpora, a fair head-to-head comparison requires evaluation sets unseen by both. We reproduce each model's exact 80/10/10 split from its fixed seed and audit, programmatically, that the evaluation rows we report contain no example present in either training or validation partition. Two such leakage-free, paired benchmarks anchor the comparison: a leftover raw Swiss set ($n=1684$, Fig.~\ref{fig:swiss}) that probes the Swiss gain, and a clean 13-group ToxiGen set ($n=424$, Fig.~\ref{fig:toxigen13}) that probes for catastrophic forgetting. In each benchmark, both models are scored on identical examples, which eliminates sampling variance from the difference.

\begin{figure}[H]
    \centering
    \includegraphics[width=\linewidth]{figures/swiss_eval.jpeg}
    \caption{Paired evaluation on the leftover raw Swiss set ($n=1684$, unseen by both models). Top: macro F1 per Swiss group and overall metrics. Bottom: per-model confusion matrices on the same texts.}
    \label{fig:swiss}
\end{figure}

\subsection{Results}
On their respective held-out test splits, ToxiGen+ is marginally stronger than ToxiGen, with both accuracy and macro F1 rising from 0.854 to 0.868. The gap widens sharply on the unseen Swiss benchmark (Fig.~\ref{fig:swiss}): macro F1 climbs from 0.898 to 0.998 and hate-class recall from 0.834 to 0.996. The confusion matrices localise the gain: of 753 truly-hate Swiss examples, the baseline misses 125 as false negatives, whereas ToxiGen+ misses only three. The largest per-group improvement is on the Portuguese community (0.842 to 0.998), the only Swiss minority without a related proxy among the original 13 identity groups. This result is particularly notable: it suggests that even a single new group, when represented by as few as 1,000 synthetic samples, is sufficient to close the detection gap for that community.

\begin{figure}[H]
    \centering
    \includegraphics[width=\linewidth]{figures/toxigen13_eval.jpeg}
    \caption{No-forgetting check on the clean 13-group ToxiGen set ($n=424$, unseen by both training loops). Top: macro F1 per original ToxiGen category, with per-model averages (dashed). Bottom: overall metrics and per-model confusion matrices.}
    \label{fig:toxigen13}
\end{figure}

On the leakage-free 13-group ToxiGen set (Fig.~\ref{fig:toxigen13}), mean per-group macro F1 rises from 0.837 to 0.848 and hate recall from 0.824 to 0.869. Several original groups improve, others regress modestly, but none collapses: the lowest post-fine-tuning per-group score is 0.72. The Swiss gains are thus acquired at no catastrophic cost to the model's coverage of the original distribution. Table~\ref{tab:results} summarises the key metrics across both benchmarks.

\begin{table}[H]
\centering
\caption{Summary of evaluation results. Macro F1 and hate-class recall for both models on the two paired benchmarks.}
\label{tab:results}
\small
\begin{tabular}{llcc}
\hline
\textbf{Benchmark} & \textbf{Metric} & \textbf{ToxiGen} & \textbf{ToxiGen+} \\
\hline
\multirow{2}{*}{Swiss ($n=1684$)} & Macro F1 & 0.898 & 0.998 \\
 & Hate recall & 0.834 & 0.996 \\
\hline
\multirow{2}{*}{ToxiGen-13 ($n=424$)} & Macro F1 & 0.837 & 0.848 \\
 & Hate recall & 0.824 & 0.869 \\
\hline
\end{tabular}
\end{table}

\subsection{Limitations}
Three caveats temper the headline numbers. First, the Swiss training and evaluation data share a single GPT-4o-mini generator, so the 0.998 macro F1 should be read as in-distribution generalisation to the target style rather than as evidence of robustness to organically-authored Swiss forum hate; external validation on annotated real-world text remains future work. Second, all results derive from a single seed-1 stratified split with neither $k$-fold cross-validation nor bootstrap confidence intervals; the headline Swiss gain is robust to this critique because it is measured on $n=1684$ paired examples, but per-group deltas on smaller subsets are noisier. Third, every metric is reported at the implicit 0.5 decision threshold (see Eq.~\ref{eq:argmax}), and a precision/recall trade-off analysis remains future work.

A fourth limitation concerns cultural bias: both \texttt{tomh/toxigen\_roberta} and \texttt{gpt-4o-mini} carry US-centric priors, so generated Swiss samples may reflect American rhetorical patterns rather than authentic Swiss online discourse. Addressing this would require native-language prompts, human annotators, and ideally a multilingual encoder.

%%%%%%%%% ADD CONCLUSION SECTION %%%%%%%%%
\section{Conclusion}
\label{sec:conclusion}
In this work, we demonstrated that \texttt{tomh/toxigen\_roberta}, despite strong performance on the original ToxiGen benchmark, fails to generalize to Swiss-specific minority groups, achieving only 26\% hate-class recall on cross-border workers, Portuguese residents, and asylum seekers. By constructing ToxiGen+, a dataset augmented with GPT-4o-mini generated samples for these groups, and fine-tuning the same RoBERTa-large architecture on it, we showed that targeted data augmentation is an effective strategy: macro F1 on the Swiss benchmark rises from 0.898 to 0.998, with no catastrophic forgetting on the original 13 ToxiGen groups.

These results highlight the importance of culturally and regionally aware dataset construction for hate speech detection. Future work could explore multilingual extensions of ToxiGen+ to cover French, German, and Italian-speaking Swiss communities, as well as the use of human annotation to validate the quality of synthetically generated samples.

% \newpage

%%%%%%%%% REFERENCES SECTION %%%%%%%%%
\bibliographystyle{ieeetr}
% This is the link to the refereces file.
\bibliography{main}

\end{document}